\section{Problemi}

\subsection{Problemi di decisione e di ottimizzazione}
La teoria della complessitá studia i problemi come oggetti matematici, indipendentemente dagli algoritmi che li risolvono. \\
La difficoltá di un problema é misurata al crescere del tempo e dello spazio di memoria necessari in funzione della dimensione dell'input.

I problemi possono essere di decisione, cioé con risposta binaria $"yes"$ o $"no"$, 
oppure di ottimizzazione, che richiedono una soluzione in output. \\

I problemi di ottimizzazione possono essere trattati come problemi di decisione, quindi
ci concentreremo su questi ultimi. Se un problema di decisione puó essere risolto in tempo polinomiale,
anche quello di ottimizzazione puó essere risolto in tempo polinomiale, tramite una ricerca binaria (a tempo logaritmico)
sul valore della soluzione ottima (vedremo esempi per $\MAXFLOW$ e $\TSP$). \\

Un problema di decisione puó essere codificato come un linguaggio $L \in \Sigma^*$ 
tale che una macchina di Turing risponde $"yes"$ a una istanza di un problema decisionale 
se la codifica dell'istanza  che riceve in input é una stringa che appartiene al linguaggio. \\

Analogamente, i problemi di ottimizzazione possono essere visti come funzioni con dominio in $\Sigma^*$
\footnote{In realtá, un problema di ottimizzazione potrebbe avere piú soluzioni, quindi anche in questo caso sarebbe meglio
parlare di relazioni. Ma per semplicitá sottintendiamo che, se esiste piú di una soluzione,
allora l'immagine della funzione sará l'insieme delle soluzioni. In ogni caso un algoritmo deterministico restituisce una sola delle tante 
soluzioni.}.


\subsection{Codifiche ragionevoli}
La stessa istanza di un problema puó essere codificata in modi diversi, con dimensioni diverse. \\

\textbf{Esempio:} \\
Se l'input è un grafo $x = (V, E)$ rappresentato tramite liste di adiacenza la dimensione dell'input è $n = |V| + |E|$. \\
Se invece il grafo viene codificato tramite matrice di adiacenza, la dimensione dell'input è $n = |V|^2$. \\
In questo caso non cambia molto, perché $|E| \leq |V|^2$. \\

\textbf{Esempio:} \\
Se l'input é un numero, viene tipicamente codificato in notazione posizionale (di solito binaria, decimale o esadecimale).
La lunghezza della rappresentazione posizionale di un numero richiede una quantità logaritmica di cifre. \\
Di conseguenza, quando l'input $x$ é un numero rappresentato in notazione posizionale, $n = log(x)$. \\
Se invece $x$ fosse rappresentato in unario, allora la lunghezza della sua codifica sarebbe uguale al numero stesso: $n = x$. \\

\subsubsection{La codifica dell'input puó alterare l'analisi di complessitá.}
\textbf{Esempio:} \\
Vogliamo calcolare il fattoriale di un numero naturale $x$. Consideriamo una versione naive che itera su tutti i numeri da 1 a x: $f(n) = f(log(x)) = x$. \\
La soluzione trovata impiega tempo esponenziale su $n$: $f(n) = O(2^n)$. \\
Si potrebbe codificare $x$ in unario: in quel caso avremmo che $x = n$ e il tempo risulterebbe essere una funzione lineare: $f(n) = O(n)$. \\

\textbf{Perché la misura che effettuiamo sia corretta, la codifica dell'input deve essere ragionevole, cioé non deve utilizzare piú spazio del necessario.}


\subsection{$\REACHABILITY$}
\label{problema:reachability}
Dato un grafo $G = (V, E)$ dove $E \subseteq V \times V$, e due vertici $s, t \in V$:
\[
\REACHABILITY = \{ (V, E, s, t) \mid s \text{ raggiunge } t \text{ in } (V, E) \} 
\]
\[
n = |V| + |E| + 2
\]
Questo problema é alla base di una importante tecnica dimostrativa, che useremo piú avanti.

\subsubsection{Algoritmo per $\REACHABILITY$}
Un algoritmo per risolvere $\REACHABILITY(G, s, t)$ consiste nell'iniziare inserendo $s$ in un insieme $S$ e marcarlo. \\
Poi iterare i seguenti passi finché $S$ non é vuoto:
\begin{enumerate}
    \item Estrarre un vertice qualsiasi $v$ da $S$,
        \subitem Per la correttezza ed efficienza di \textit{questo} algoritmo, non importa quale vertice estrae. 
        Vedremo che, per la riduzione di $\MAXFLOW$ a $\REACHABILITY$ 
        (\ref{riduzione:maxflow_reachability}), 
        conviene estrarre il vertice piú vicino a $s$,
        implementando $S$ come una coda (Breadth-First Search).
        Questo garantisce di trovare sempre il cammino piú corto da $s$ a $t$.
    \item Inserire in $S$ tutti i vertici adiacenti a $v$ non marcati, e marcarli.
\end{enumerate}
Alla fine, se $t$ é marcato, allora $s$ raggiunge $t$. \\
Nel caso peggiore esplora tutti gli archi del grafo una sola volta, quindi l'algoritmo é polinomiale: $O(|E|) = O(|V|^2) = O(n^2)$. \\

Un altro modo per risolvere il problema é mostrato nella dimostrazione del teorema di Savitch (\ref{theorem:savitch}). \\


\newpage
\subsection{$\MAXFLOW$}
Sia $N = (V, E, s, t, c)$ un \textbf{network}, cioé:
\begin{itemize}
	\item un grafo $(V, E)$, 
	\item la funzione di \textbf{capacitá} $c: E \rightarrow \mathbb{N}$,
	\begin{itemize}
        \item $c_{max} = \max_{(u, v) \in E} c(u, v)$
        \item supponiamo che i valori di $c$ siano rappresentati in binario. \\
        $|c| = |E| \cdot \log_2(c_{max})$
    \end{itemize}
	\item un vertice sorgente $s \in V$ senza archi entranti 
	\item un vertice pozzo $t \in V$ senza archi uscenti.
\end{itemize} 
Un \textbf{flusso} $f$ in $N$ é una funzione $f: E \rightarrow \mathbb{N}$ tale che:
\begin{itemize}
    \item $\forall (u, v) \in E : f(u, v) \leq c(u, v)$. 
        \subitem Il flusso su un arco non puó superare la sua capacitá.
    \item $\forall v \in V \setminus \{s, t\} : \sum_{(u, v) \in E} f(u, v) = \sum_{(v, w) \in E} f(v, w)$. 
        \subitem La quantitá di flusso che entra in un vertice é uguale a quella che esce.
\end{itemize}
Il \textbf{valore} di un flusso é definito come la somma degli archi uscenti da $s$: $\sum_{(s, v) \in E} f(s, v)$. \\
O, equivalentemente, come la somma degli archi entranti in $t$: $\sum_{(u, t) \in E} f(u, t)$. \\

$\MAXFLOW$ é il problema di trovare un flusso $f$ in un network $N$ tale che il suo valore sia massimo.
É un problema di ottimizzazione, rappresentato da una funzione 
\[
\MAXFLOW: (V, E, s, t, c) \in \Sigma^* \rightarrow (f: E \rightarrow \mathbb{N}) \in \mathbb{N}^E
\]
\[
n = |V| + |E| + 2 + |c|
\]

\subsubsection{$\MAXFLOW_D$}
Consideriamo il problema decisionale che chiede se esiste un flusso di valore almeno $k$ in un network $N$.
\[
\MAXFLOW_D = \{ (V, E, s, t, c, k) \in \Sigma^* \mid (V, E, s, t, c) \text{ ha un flusso di valore} \geq k \}
\]
\[
n = |V| + |E| + 2 + |c| + \log_2(k)
\]

\subsubsection{Da $\MAXFLOW_D$ a $\MAXFLOW$}
Possiamo risolvere il problema di ottimizzazione $\MAXFLOW$ utilizzando una soluzione per la sua versione decisionale $\MAXFLOW_D$: \\
Consideriamo il massimo flusso possibile, dato dal minimo tra la somma delle capacitá uscenti da $s$ e la somma delle capacitá entranti in $t$. \\
\[
m = \min\left(\sum_{v : (s,v) \in E} c(s, v), \sum_{u : (u, t) \in E} c(u, t)\right)
\]
Eseguiamo una ricerca binaria tra $0$ e $m$, interrogando $\MAXFLOW_D$ per trovare il valore massimo del flusso. \\
Eseguiamo quindi al piú $\log_2(m)$ interrogazioni. \\
$m \leq |E| \cdot c_{max} = O(2^n)$. \\
La codifica dei valori di $c$ é logaritmica, di conseguenza $c_{max}$ é esponenziale rispetto alla dimensione dell'input. \\
$O(\log_2(m)) = O(\log_2(|E| \cdot c_{max})) = O(\log_2(2^n)) = O(n)$. \\
Il numero di interrogazioni é polinomiale rispetto alla dimensione dell'input. \\
Quindi una soluzione polinomiale per $\MAXFLOW_D$ implica una soluzione polinomiale per $\MAXFLOW$.


% \newpage
\subsection{$\BIPARTITEMATCHING$}
Sia $G = (U, V, E)$ un grafv vio bipartito. $U, V$ sono due insiemi di vertici con la stessa cardinalitá $|U| = |V|$. \\ 
Gli archi vanno solo da $U$ a $V$: $E \subseteq U \times V$. \\
Un \textbf{matching} è un insieme di archi $M \subseteq E$ tale che ogni vertice di $U$ ha un arco in $M$ verso uno e un solo vertice di $V$.
\[
\BIPARTITEMATCHING = \{ (U, V, E) \mid (U, V, E) \text{ ha un matching} \}
\]


% \newpage
\subsection{$\SAT$}
\label{problema:sat}
Se $PL$ é l'insieme di tutte le formule della logica proposizionale, e $AP$ é l'insieme di tutte le proposizioni atomiche:
\[
\SAT = \{ \phi \in PL \mid \exists M \subseteq AP : M \models \phi \} 
\]
\[
n = |\phi|
\]

É un problema che richiede di esplorare uno spazio esponenziale di possibili assegnazioni a ogni proporizione atomica in $\phi$. \\
Un problema affine é quello del \textbf{model checking}: dato un modello $M$, verifica se $M \models \phi$.
Questo problema é molto piú semplice da risolvere, infatti richiede tempo polinomiale. \\
Un algoritmo non deterministico per $\SAT$ puó essere visto come un algoritmo di model checking, 
dove il modello $M$ é scelto in maniera non deterministica. \\
Basterebbe infatti rappresentare $\phi$ tramite l'albero sintattico ed effettuare una visita in post-order,
calcolando il valore di veritá dei nodi interni e restituendo il valore del nodo radice. \\
Il valore di veritá delle foglie, cioé delle proposizioni atomiche, é scelto in maniera non deterministica. 
Questo vuol dire che l'algoritmo controlla in maniera simultanea tutti i modelli possibili. \\

\subsubsection{$\CNFSAT$}
Soddisfaciilibtá di una formula in forma normale congiuntiva (CNF), cioé congiunzione di clausole disgiuntive (forma $\bigwedge(\bigvee(...))$). \\

\subsubsection{$\kSAT$}
Soddisfaciilibtá di una formula in forma normale congiuntiva (CNF), dove ogni clausola ha esattamente $k$ letterali. \\


\subsection{$\HAMILTONPATH$}
\label{problema:hamiltonpath}
Sia $G = (V, E)$ un grafo.

Un ciclo hamiltoniano é un ciclo che passa per tutti i vertici del grafo una sola volta.
Puó essere rappresentato da una funzione biettiva $\pi$ che codifica una permutazione dei vertici del grafo, 
associando ogni vertice a una e una sola posizione nel ciclo hamiltoniano.
Inoltre deve esistere un arco tra $\pi(|V|)$ e $\pi(1)$.\\
É un problema di ottimizzazione, rappresentato da una funzione:
\[
\HAMILTONPATH: (V, E) \in \Sigma^* \rightarrow (\pi : \{1,\dots,|V|\} \rightarrow V) \in V^{\{1,\dots,|V|\}} :
\]
\[
\pi \text{ biettiva} \quad \land \quad \forall i \in \{1,\dots,|V|\} (\pi(i), \pi((i+1) \bmod {|V|})) \in E
\]
La dimensione dell'input é \\
\[
n = |V| + |E|
\]


\subsection{$\TSP$}
\label{problema:tsp}
Sia $G = (V, E, w)$ un grafo pesato, dove $w: E \rightarrow \mathbb{N}$ é la funzione di peso degli archi. \\
Supponiamo che i valori di $w$ siano rappresentati in binario: \\
$|w| = \log_2(w_{max})$ \\
$w_{max} = \max_{(u, v) \in E} w(u, v)$. \\

Chiamiamo $w(\pi)$ il costo di un ciclo hamiltoniano $\pi$:
\[
w(\pi) = \sum_{i = 1}^{|V|} w(\pi(i), \pi((i+1) \bmod {|V|}))
\]

Il problema del commesso viaggiatore, o Traveling Salesman Problem, consiste nel trovare
un ciclo hamiltoniano di costo minimo in un grafo pesato. \\
Intuitivamente, immaginiamo un commesso viaggiatore che vuole visitare tutte le cittá della regione una sola volta, e tornare alla cittá di partenza nel minor tempo possibile. \\
Anche questo é un problema di ottimizzazione, rappresentato da una funzione:
\[
\TSP: (V, E, w) \in \Sigma^* \rightarrow {\arg\min}_{\pi \in \HAMILTONPATH(V, E)} w(\pi)
\]
\[
n = |V| + |E| + |w|
\]


\subsubsection{$\TSP_D$}
Dati $G = (V, E, w)$ e un numero naturale $B \in \mathbb{N}$ detto \textit{budget}:
\[
\TSP_D = \{ (V, E, w, B) \mid (V, E) \text{ ha un ciclo hamiltoniano di costo } \leq B \}
\]
\[
n = |V| + |E| + |w| + \log_2(B)
\]

\subsubsection{Da $\TSP_D$ a $\TSP$}
Come per $\MAXFLOW$, anche in questo caso possiamo risolvere il problema di ottimizzazione $\TSP$,
interrogando $\TSP_D$ per trovare un budget minimo per cui la risposta é $"yes"$. \\
Il budget minimo si cerca tramite una ricerca binaria (logaritmica) su tutti i possibili costi dei cicli hamiltoniani. \\

Siccome i valori di $w$ sono rappresentati in notazione posizionale, e $|w| < n$ allora il valore massimo
del peso di un arco é al piú $2^n$. \\
Moltiplicandolo per il numero massimo di archi in un ciclo hamiltoniano, che é $|V| \leq n$,
otteniamo che il costo massimo di un ciclo hamiltoniano é al piú $n \cdot 2^n$. \\
Una ricerca binaria (logaritmica) su $n \cdot 2^n$ possibili valori richiede al piú $\log(n) + n$ interrogazioni. \\

Dopo averlo trovato, fissiamo il budget $B$ e continuiamo a interrogare $\TSP_D$, modificando il grafo: \\ 
a ogni interrogazione, eliminiamo un arco (per esempio settando il suo costo a $B + 1$). \\
Se la risposta é sempre $"yes"$, allora l'arco non fa parte del ciclo hamiltoniano ottimo, e possiamo escluderlo anche dalle interrogazioni successive. \\
Altrimenti l'arco fa parte del ciclo hamiltoniano ottimo, e bisogna ripristinarlo. \\

Dopo aver interrogato tutti gli archi, in al piú $n^2$ interrogazioni ($|E| \leq |V|^2 \leq n^2$), 
gli archi che non sono stati eliminati sono quelli che formano il ciclo hamiltoniano ottimo. \\
Infatti se un arco rimane nell'insieme finale  ma non appartiene al ciclo hamiltoniano, 
allora nel momento in cui è stato testato per la rimozione, la risposta sarebbe dovuta essere "yes" 
(perché il ciclo non lo usa ed è ancora disponibile in quel momento), contraddicendo il fatto che 
sia stato ripristinato. \\

Quindi, possiamo risolvere $\TSP$ invocando $\TSP_D$ un numero polinomiale di volte: $O(n + n^2) = O(n^2)$.

